\section{总结与延伸}

本讲从“为什么机器人还不会收拾房间”出发，最终落在一个清晰的系统分解：foundation model 提供 semantic prior，robot data 提供 action grounding，interface 决定模型输出的抽象层级，controller 处理 dynamics 与 constraint，evaluation 检查 transfer，safety layer 限制物理后果。PaLM-E 证明 sensor embedding 与 language model 可以共同训练；RT-2 把 action token 接进 VLM，形成 VLA；RT-X 测试 cross-embodiment positive transfer；Language to Rewards 则让 LLM 写 objective、MPC 写 trajectory。

\subsection{六条可迁移的工程启示}

\begin{enumerate}
  \item \textbf{先定义接口，再选模型。} 决定输出是 skill、action token、trajectory、reward 还是 constraint，比盲目增大参数更基础。
  \item \textbf{把 data mixture 当成设计对象。} Web data、robot success、failure、recovery 与多 embodiment 数据需要不同权重和审计。
  \item \textbf{分开 semantic generalization 与 motor generalization。} 会识别新对象不等于会产生新接触动作；评价指标必须拆开。
  \item \textbf{用 ablation 证明 transfer 路径。} 只看大模型更好，无法区分 scale、co-fine-tuning、robot data 与 web prior 的作用。
  \item \textbf{保留 controller 与 safety layer 的可见性。} 即使 policy 端到端，执行层仍应有 rate/force/geometry guardrail 与 emergency stop。
  \item \textbf{考虑 hybrid bootstrap。} Reward optimizer 可生成 diverse data，VLA 可 distill 成快速 policy，两者可以形成 data flywheel。
\end{enumerate}

\begin{importantbox}{把本讲压缩成一句话}
Foundation models 对机器人最重要的贡献，不是让语言模型直接控制一切，而是让我们重新设计\textbf{共享表示、训练数据与任务接口}；真正可用的 embodied system 仍必须把语义、动作、动力学、评价和安全连接成闭环。
\end{importantbox}

\subsection{最容易误读的三句话}

\begin{warningbox}{不要把局部证据扩大为通用结论}
“RT-2 有 emergent skills”不等于学会任意新动作；“RT-X 有 positive transfer”不等于 robotics scaling law 已成熟；“LLM 能写 reward”不等于生成代码可直接上真实机器人。三句话都需要附带 action support、dataset boundary、controller、validation 与 safety 条件。
\end{warningbox}

\subsection{拓展阅读}

\begin{itemize}
  \item Driess et al., \href{https://arxiv.org/abs/2303.03378}{\emph{PaLM-E: An Embodied Multimodal Language Model}}
  \item Brohan et al., \href{https://arxiv.org/abs/2212.06817}{\emph{RT-1: Robotics Transformer for Real-World Control at Scale}}
  \item Brohan et al., \href{https://arxiv.org/abs/2307.15818}{\emph{RT-2: Vision-Language-Action Models Transfer Web Knowledge to Robotic Control}}
  \item Open X-Embodiment Collaboration, \href{https://arxiv.org/abs/2310.08864}{\emph{Open X-Embodiment: Robotic Learning Datasets and RT-X Models}}
  \item Yu et al., \href{https://arxiv.org/abs/2306.08647}{\emph{Language to Rewards for Robotic Skill Synthesis}}
  \item Lynch et al., \href{https://arxiv.org/abs/2206.04717}{\emph{Interactive Language: Talking to Robots in Real Time}}
\end{itemize}

阅读这些论文时可以沿用本讲的五格检查表：observation、model、interface、controller、evidence。只要其中一格被“端到端”或“emergent”一词遮住，就继续追问数据来自哪里、动作怎样表达、物理约束由谁处理、失败如何被记录，以及什么实验能够否定当前解释。

\end{document}
